\chapter{Apraxia of Speech}

\section*{Definition and Overview}
Apraxia of speech is a neurological disorder in which the brain has difficulty planning and programming the precise movements needed to produce speech, even though the muscles of the lips, tongue, jaw, and palate are structurally intact and strong. The person knows the words they wish to say and understands language, but the brain cannot reliably generate the sequence of motor commands required to make the sounds in the correct order. The disorder can appear in childhood (childhood apraxia of speech) or be acquired in adulthood after damage to the brain, most commonly a stroke. It is a disorder of motor planning rather than of language comprehension or muscle weakness, which distinguishes it from aphasia and dysarthria.

\section*{Epidemiology}
Childhood apraxia of speech is relatively rare, affecting roughly one to two children in every thousand, and it is more common in boys. Acquired apraxia of speech in adults most often follows a stroke affecting the left cerebral hemisphere, but it can also result from head injury, brain surgery, or tumour. It may develop in the course of certain neurodegenerative conditions, such as progressive apraxia of speech, corticobasal syndrome, and some forms of dementia. Apraxia of speech frequently coexists with aphasia, a disorder of language itself, which can make the two difficulties hard to separate clinically.

\section*{Aetiology and Pathophysiology}
Apraxia of speech is caused by damage or underdevelopment of the brain regions that plan and programme the movements of the articulators for speech, particularly the left frontal and parietal areas that coordinate the lips, tongue, jaw, velum, and larynx.

\begin{itemize}
    \item \textbf{Childhood apraxia:} the cause is often unknown; genetic factors may play a part, and it can occur alongside other developmental conditions such as autism spectrum disorder or epilepsy
    \item \textbf{Stroke:} the commonest cause in adults, especially when it involves the language and motor-planning areas of the dominant hemisphere
    \item \textbf{Head injury:} traumatic damage to the speech-planning regions of the brain
    \item \textbf{Brain tumour or surgery:} injury to speech-related cortex during treatment
    \item \textbf{Degenerative conditions:} progressive loss of brain tissue, as in progressive apraxia of speech or corticobasal syndrome
    \item \textbf{Other neurological disease:} some forms of dementia or motor neurone disease can involve speech apraxia
\end{itemize}

\section*{Clinical Features}
In both children and adults, the core feature is inconsistent, effortful speech that typically worsens when the person attempts longer or more complex utterances.

\begin{itemize}
    \item \textbf{Childhood apraxia:} late first words, a restricted sound repertoire, and limited vocabulary for age
    \item Inconsistent errors, with the same word produced differently on repeated attempts
    \item Distorted sounds and particular difficulty sequencing sounds within longer or complex words
    \item Visible groping or searching movements of the mouth when trying to speak
    \item Automatic speech, such as greetings or well-rehearsed phrases, may be better than deliberate speech
    \item Comprehension of language is usually far better than speech output
    \item In adults: slow, halting, effortful speech with pauses, sound distortions, and difficulty initiating words
    \item Frustration, and in children behavioural difficulties related to the communication problem
\end{itemize}

\section*{Differential Diagnosis}
Several other speech and language disorders resemble apraxia of speech and must be distinguished, because their management differs.

\begin{itemize}
    \item Dysarthria, in which the speech muscles are weak, spastic, or poorly coordinated
    \item Aphasia, a disorder of language itself rather than of speech motor planning
    \item Phonological disorder in children, where the sound system of the language is learned incorrectly
    \item Hearing impairment, which limits the development or control of speech
    \item Stuttering or cluttering, which involve fluency rather than motor planning
    \item General developmental delay or intellectual disability
    \item Voice disorders or articulation errors due to dental and mouth problems
\end{itemize}

\section*{When to Seek Medical Care}
See a doctor if a child is not meeting expected speech milestones, has little or no babbling by the age of 12 months, or is losing previously acquired words. Early referral to a speech pathologist offers the best chance of improvement. In adults, any new or worsening difficulty with speech should be assessed promptly, because it may indicate an underlying neurological cause.

\textbf{Call 000 (or local emergency services) for:}
\begin{itemize}
    \item Sudden difficulty speaking, especially with weakness of one side of the body, facial droop, or confusion, which may indicate a stroke
    \item Loss of consciousness or a seizure
    \item Sudden difficulty breathing or swallowing
\end{itemize}

\section*{Diagnosis and Investigations}
There is no single test for apraxia of speech. The diagnosis is made by a speech pathologist on the basis of a careful assessment of motor speech, supported by investigations into any underlying cause.

\begin{itemize}
    \item \textbf{Speech and language assessment:} observation of repetition, naming, and connected speech to identify the characteristic pattern of errors
    \item \textbf{Developmental history:} for children, a review of speech milestones, hearing, and other aspects of development
    \item \textbf{Hearing test:} audiometry to exclude hearing loss as a cause of poor speech
    \item \textbf{Brain imaging:} CT or MRI in adults to identify stroke, tumour, or other structural damage
    \item \textbf{Neurological assessment:} examination for other features of brain or movement disorders
    \item \textbf{Genetic testing:} sometimes offered for children in whom an inherited form is suspected
    \item \textbf{Other assessments:} of receptive and expressive language, cognition, and functional communication needs
\end{itemize}

\section*{Management}
Treatment is led by a speech pathologist and focuses on retraining the brain to plan and sequence speech movements accurately.

\begin{itemize}
    \item \textbf{Speech therapy:} intensive, frequent sessions using repetitive practice, modelling, and multisensory cues that allow the person to see, hear, and feel the target sounds
    \item \textbf{Home practice:} short, frequent practice tasks recommended by the therapist between sessions
    \item \textbf{Augmentative and alternative communication (AAC):} signing, picture boards, or speech-generating devices where speech is severely limited
    \item \textbf{Children:} early intervention, parent coaching, and support within kindergarten and school
    \item \textbf{Adults:} rehabilitation after stroke or head injury, including treatment of any coexisting aphasia
    \item \textbf{Psychological support:} help with frustration, confidence, and social participation for the person and family
    \item \textbf{Treating the cause:} management of any underlying neurological condition, such as the secondary prevention of stroke
\end{itemize}

\section*{Complications}
\begin{itemize}
    \item Ongoing communication difficulty and frustration for the person and family
    \item In children: delayed literacy and learning, teasing, and low self-esteem
    \item Behavioural problems linked to the inability to communicate effectively
    \item Social isolation for both children and adults
    \item In adults: reduced independence, employment prospects, and quality of life
    \item Anxiety and depression related to the loss of communication
\end{itemize}

\section*{Prognosis and Outcomes}
With intensive, consistent therapy, many children with apraxia of speech make substantial gains and go on to communicate effectively, although some continue to have subtle speech difficulties into adolescence and adulthood. Outcomes depend on how early treatment begins and on its intensity and consistency. In adults, recovery is highly variable: many improve considerably with rehabilitation after stroke, but recovery may be incomplete, and in progressive neurological conditions the speech difficulty tends to worsen over time. Early recognition, skilled therapy, and good family support are the strongest determinants of a favourable outcome.

\section*{Prevention and Self-Care}
Apraxia of speech itself cannot be prevented, but early recognition and treatment significantly improve outcomes. Families can do much to support communication in daily life.

\begin{itemize}
    \item Seek a speech pathology assessment early if speech milestones are delayed
    \item Practise speech exercises daily under the therapist's guidance
    \item Be patient, and give the person extra time to speak without pressure
    \item Use gestures, pictures, or writing to support understanding and reduce frustration
    \item Reduce background noise and gain the person's attention before speaking
    \item Reduce stroke risk in adults through blood pressure control, not smoking, and a healthy lifestyle
    \item Seek support groups and counselling for the person and the family
\end{itemize}

\section*{Sources and Further Reading}
The following organisations provide reliable information on apraxia of speech for patients, families, and healthcare students.

\begin{itemize}
    \item NHS. \textit{Speech and language therapy: apraxia of speech information.} https://www.nhs.uk/conditions/speech-and-language-therapy/
    \item Mayo Clinic. \textit{Childhood apraxia of speech.} https://www.mayoclinic.org/diseases-conditions/childhood-apraxia-of-speech/
    \item MedlinePlus. \textit{Speech and communication disorders.} https://medlineplus.gov/speechandcommunicationdisorders.html
    \item National Institute on Deafness and Other Communication Disorders (NIDCD). \textit{Apraxia of speech.} https://www.nidcd.nih.gov/health/apraxia-speech
    \item American Speech-Language-Hearing Association (ASHA). \textit{Childhood apraxia of speech.} https://www.asha.org/public/speech/disorders/childhood-apraxia-of-speech/
    \item Stroke Foundation. \textit{Communication difficulties after stroke.} https://strokefoundation.org.au/what-we-do/support-for-you/communication-after-stroke
\end{itemize}
